% TOP_PROBLEM: {"id": 30006977, "title": "Global Gan–Gross–Prasad Conjecture (Fixed-Representation Form)", "category_id": 1, "category": {"id": 1, "name": "number_theory", "display_name": "Number Theory", "description": "Properties of integers, prime numbers, Diophantine equations.", "slug": "number-theory", "order_index": 1, "created_at": "2026-07-31T15:26:25.670Z"}, "status": "open"}
% FIRST_POSTED: null
% !TeX program = lualatex
% Compile twice: lualatex 30006977.tex
% All references and prose are embedded; no BibTeX database or external inputs.
\documentclass[11pt,a4paper]{article}
\usepackage[margin=24mm,headheight=14pt]{geometry}
\usepackage{fontspec}
\setmainfont{Latin Modern Roman}
\setsansfont{Latin Modern Sans}
\setmonofont{Latin Modern Mono}
\usepackage{amsmath,amssymb,mathtools}
\usepackage{unicode-math}
\setmathfont{Latin Modern Math}
\usepackage{microtype}
\usepackage{enumitem,xurl,needspace}
\usepackage[dvipsnames]{xcolor}
\usepackage{hyperref}
\hypersetup{unicode=true,colorlinks=true,linkcolor=MidnightBlue,urlcolor=MidnightBlue,
 pdftitle={Research attempt: Problem 30006977},pdfauthor={Research notebook},bookmarksnumbered=true}
\usepackage{bookmark}
\usepackage{tocloft}
\setlength{\cftsecnumwidth}{3em}
\cftsetpnumwidth{2.5em}
\cftsetrmarg{4em}
\usepackage{titlesec}
\titleformat{\section}{\Large\bfseries\raggedright}{\thesection}{0.7em}{}
\usepackage{fancyhdr}
\pagestyle{fancy}\fancyhf{}
\fancyhead[L]{\small\sffamily Open Problems Math | Research notebook}
\fancyhead[R]{\small Problem 30006977 | 27 September 2026}
\fancyfoot[C]{\thepage}
\setlength{\parindent}{0pt}
\setlength{\parskip}{3pt plus 1pt}
\setlength{\emergencystretch}{3em}
\setcounter{tocdepth}{1}
\setcounter{secnumdepth}{1}
\setlist[itemize]{leftmargin=3em,itemsep=5pt,topsep=4pt,parsep=0pt}
\allowdisplaybreaks
\newcommand{\N}{\mathbb{N}}
\newcommand{\Z}{\mathbb{Z}}
\newcommand{\Q}{\mathbb{Q}}
\newcommand{\R}{\mathbb{R}}
\newcommand{\C}{\mathbb{C}}
\newcommand{\F}{\mathbb{F}}
\newcommand{\A}{\mathbb{A}}
\newcommand{\PP}{\mathbb P}
\newcommand{\E}{\mathbb{E}}
\newcommand{\eps}{\varepsilon}
\newcommand{\dd}{\,\mathrm d}
\newcommand{\sref}[2]{\hyperlink{source:#1:#2}{[S#2]}}
\newcommand{\eref}[2]{\hyperlink{extra:#1:#2}{[E#2]}}
% Turn the next switch off for a shorter reading copy. The quotations preserve
% every exactTarget field, including qualifications omitted from a short summary.
\newif\ifshowcatalogue
\showcataloguetrue
\newcommand{\cataloguescope}[1]{\ifshowcatalogue
 \paragraph{Exact scope in the supplied catalogue.}
 \begin{quote}\small #1\end{quote}\fi}
\usepackage{etoolbox}
\AtBeginEnvironment{itemize}{\small}
\numberwithin{equation}{section}
% Standalone export. Original research content preserved.
\begin{document}
\setcounter{section}{0}
% ================================================================
% problemId: 30006977
% Canonical problem ID: 30006977
% exactTarget SHA-256: 311d42f410b25a3abd28d1e18da7ee300644f35f867c066e172d96ea66d496b7
\clearpage
\section*{\texorpdfstring{Global Gan–Gross–Prasad Conjecture (Fixed-Representation Form)}{Global Gan–Gross–Prasad Conjecture (Fixed-Representation Form)}}\label{30006977}
\subsection{Definitions and mathematical statement}
For the relevant classical-group pair, let $H$ be the subgroup in Gan--Gross--Prasad, $\nu=\bigotimes'_v\nu_v$ its Bessel or Fourier--Jacobi data, and $R$ the prescribed representation defining the finite $L$-function. For each fixed tempered cuspidal automorphic representation $\pi$ occurring with multiplicity one, the assertion is
\[
 \mathcal F(\nu)|_\pi\ne0
 \quad\Longleftrightarrow\quad
 \left(\forall v,\ \operatorname{Hom}_{H(F_v)}(\pi_v,\nu_v)\ne0\right)
 \ \text{and}\ L(\pi,R,\tfrac12)\ne0.
\]
The left side says that the source-defined global period is not identically zero on $\pi$. The local Hom-spaces test distinction at every place. The form signs, covers, auxiliary characters, and inverse twist in $R$ are retained from Sections 23--24. Neither a packet-existence statement nor a refined numerical period formula is substituted.
\par\smallskip\textit{Formulation and concept references:} \sref{30006977}{1}.
\cataloguescope{For every number field F, let E be F, a quadratic extension, or F\ensuremath{\times}F, and let W\ensuremath{\subset}V be a nondegenerate relevant pair with form sign epsilon, split W-perp, and epsilon*(-1)\textasciicircum{}dim(W-perp)=-1; a two-dimensional orthogonal V or W is required to be nonsplit. Use exactly the groups, covers, Bessel/Fourier–Jacobi representation nu and auxiliary data of Gan–Gross–Prasad Section 23, and R of Section 24, including the skew-hermitian character mu and its matching inverse twist in R. For every irreducible tempered representation pi occurring with multiplicity one in the cuspidal automorphic spectrum, Conjecture 24.1 asserts: the global period F(nu) is nonzero on pi if and only if Hom\_H(F\_v)(pi\_v,nu\_v) is nonzero for every place v and L(pi,R,1/2) is nonzero. Quantification is over all data permitted there, including relevant orthogonal, hermitian, symplectic/metaplectic and skew-hermitian cases. This is the fixed-representation equivalence, not the refined period identity or packet-existence formulation.}
\subsection{Short English statement}
For each allowed automorphic representation, is its global period nonzero exactly when every local period is allowed and the central $L$-value is nonzero?
% BEGIN RESEARCH ADDITION 138
\subsection{Research attempt}

\paragraph{Source audit and current frontier.}
Sections 23--24 of Gan--Gross--Prasad were inspected, retaining the conjugate
Weil representation in the Fourier--Jacobi pairing and the $\mu^{-1}$ twist
in $R$. Conjecture 24.1 is for the fixed representation, not merely existence
somewhere in a packet. Its split corank-one Rankin--Selberg remarks are
consistency checks, not substitutes for all allowed cases. \sref{30006977}{1}
Boisseau--Lu--Xue's 2026 manuscript, the third in their unitary
Fourier--Jacobi series, supplies a global nonvanishing statement, local
detection and a refined formula in Theorems 1.1, 1.3 and 1.4. This is substantial
specific progress, not a result for all orthogonal and symplectic/metaplectic
cases in the catalogue. Those are the cited manuscript's theorems; its full
proofs have not independently been verified here. \eref{30006977}{1}

\paragraph{Attack route.}
Local multiplicity one suggests factorizing the global period and determining
its scalar by a relative trace formula. A support comparison, identifying
only whether the scalar vanishes, would already suffice. Alternatively one
might globalize local nonzero test vectors. We show exactly why the latter
stops at a scalar, and prove a spectral-filter identity exposing the tail
estimate required to isolate one representation rather than a packet sum.

\paragraph{Attempt: local uniqueness leaves a global scalar.}
Let $E_v$ be the algebraic smooth local test space: $\pi_v$ in the Bessel
case, or $\pi_v\otimes\overline{\nu_v}$ when the Fourier--Jacobi pairing
is packaged as an invariant linear form. This notation retains the source's
dual conventions. Assume for the following lemma that the relevant invariant
form spaces have dimension one, with nonzero forms $\ell_v$, and that
reference vectors can be normalized by $\ell_v(e_v)=1$, compatibly with the
restricted tensor product outside finitely many places. These local and
unramified hypotheses are not proved anew for all catalogue cases.
\sref{30006977}{1}

For any invariant linear form $P$ on $E=\bigotimes'_vE_v$,
\begin{equation}\label{30006977:scalar}
 P(\otimes_v f_v)=c\prod_v\ell_v(f_v),\qquad c=P(\otimes_v e_v).
\end{equation}
The product is finite, since $f_v=e_v$ almost everywhere. Freeze all but
one factor: invariance makes the resulting linear form a multiple of
$\ell_v$, and evaluation at $e_v$ determines the multiple. Induction over
the finitely many modified factors proves \eqref{30006977:scalar}; linearity
extends it to the algebraic tensor product. If a local form space is zero,
the same argument kills every pure tensor. On completed smooth spaces,
continuity and density are needed to extend that conclusion. Convergence
and continuity of the actual cuspidal period must remain part of the
source-defined construction. \sref{30006977}{1}

Thus a nonzero global period forces the local conditions. Conversely all
local conditions construct an abstract nonzero tensor functional, but do
not show that the \emph{automorphic} period has $c\ne0$. The zero form
satisfies the same invariance identities. Local uniqueness therefore cannot
supply arithmetic nonvanishing. This argument also does not prove the
necessary central-$L$ condition.

\paragraph{An explicit spectral-localization requirement.}
Let $e_0,e_1,\ldots$ be orthonormal eigenvectors of a commuting family of
diagonal normal operators, and let $b_j\ge0$ be spectral weights. Suppose
for each $1\le i\le M$ an operator $T_i$ has distinct eigenvalues
$t_i(0),t_i(i)$. Define
\[
 Q_M=\prod_{i=1}^M\frac{T_i-t_i(i)I}{t_i(0)-t_i(i)},
 \qquad Q_Me_j=q_M(j)e_j.
\]
Then $q_M(0)=1$ and $q_M(i)=0$ for $1\le i\le M$. Consequently, whenever
the series converges,
\begin{equation}\label{30006977:filter}
 \sum_{j\ge0}b_j|q_M(j)|^2
       =b_0+\sum_{j>M}b_j|q_M(j)|^2.
\end{equation}
A comparison calculation with a positive lower limit on the left side and
\[
 \sum_{j>M}b_j|q_M(j)|^2\longrightarrow0
\]
would force $b_0>0$. In a correctly regularized relative trace formula the
intended $b_0$ detects the period of the fixed representation, and the
comparison side would detect its central $L$-value. This is an explicit
attack requirement, not a trace formula claimed for the unsolved cases.
Isolation and regularization are substantive steps of the contemporary
unitary argument. \eref{30006977}{1}

Finite interpolation alone does not prove the tail condition. A factor
separating eigenvalues $0$ and $\epsilon$ is $1-T/\epsilon$, whose magnitude
at eigenvalue $1$ is $|1-1/\epsilon|$. Increasingly delicate spectral
separation can amplify the remaining spectrum without bound. If two
representations have the same eigenvalues for every chosen Hecke operator,
no such filter separates them at all. Positivity of a packet sum then gives
existence of some nonzero summand, not positivity of the prescribed one.

\paragraph{Why an unnormalized Euler product does not determine the scalar.}
Nonzero local factors can have zero limiting product: exactly
$\prod_{j=2}^m(1-1/j)=1/m$. In \eqref{30006977:scalar} the unramified factors
were normalized to one; the remaining $c$ is genuinely undetermined, not
made nonzero by a convergent-product argument. Its zero set can be identified
with that of $L(\pi,R,1/2)$ only through global input. A functional equation
with sign $+1$ also does not force a nonzero central value: $(s-1/2)^2$ is an
even function with a central zero. These examples test the logic, not
construct automorphic counterexamples.

\paragraph{Stress test.}
The tensor calculation explicitly assumes one-dimensional form spaces and
compatible reference vectors, and does not identify local/global measures
by fiat. The filter identity assumes a discrete convergent sum; continuous
or residual spectra need separate control. Unramified Hecke data may only
isolate a packet. Finally, dropping the inverse $\mu$ twist changes the
$L$-function itself. None of those points is discharged by saying
``multiplicity one.'' The finite script checks the scalar and filter
identities in small models, not an automorphic trace formula.

\paragraph{Outcome.}
\textbf{Outcome: Exploratory approach with unresolved gap.}
The attempt proves an algebraic scalar reduction and a precise spectral
isolation identity, while rejecting local-test-vector nonvanishing,
packet-sum positivity and nonzero-factor multiplication as substitutes for
the missing theorem. A global support comparison with controlled remainder
for every allowed group and fixed representation is still required. No
extension of the unitary results to the full target is claimed.

% END RESEARCH ADDITION 138
\subsection{Sources}
\begin{itemize}\raggedright
\item[\textnormal{[S1]}]\hypertarget{source:30006977:1}{} Gan, Gross and Prasad, Symplectic local root numbers, central critical L-values, and restriction problems in the representation theory of classical groups, §§23–24, Conjecture 24.1.
\url{https://arxiv.org/pdf/0909.2999}.
{\small\textit{Catalogue formulation source.}}
% BEGIN RESEARCH REFERENCES 138
\item[\textnormal{[E1]}]\hypertarget{extra:30006977:1}{} P. Boisseau, W. Lu and H. Xue,
\emph{The global Gan--Gross--Prasad conjecture for Fourier--Jacobi periods
on unitary groups III: Proof of the main theorems}, arXiv:2601.01738v1
(2026), Theorems 1.1, 1.3, 1.4 and Sections 1.4--1.5.
\url{https://arxiv.org/html/2601.01738v1}.

% END RESEARCH REFERENCES 138
\end{itemize}

\end{document}
